\documentclass{article}
\usepackage{pathfinder-readable}

\title{What an Occluded Handover Can Tell Its Receiver}

\begin{document}
\maketitle

\begin{abstract}
Paper Q studies how reports and rewards can direct an agent with private information, while paper P
asks whether a donor's wrist can locate an object hidden during handover. The thread looked for a
limit on such control when wrist and motor observations leave the object's orientation unclear. The
peers derived conditional bounds on a receiver's success and on rewards based on coarse
measurements, and checked a failure of one of Q's reward menus under a changed signal assumption.
They also found a repair for a stipulated signal channel. The thread paused because P supplies
neither measured ambiguous histories nor tests showing that the hidden orientations require
different receiving actions.
\end{abstract}

\section{Introduction}

Q, \emph{Mechanism Design for Alignment and Control}, studies an agent whose preferences, possible
actions and information may be private. A human can ask for reports, recommend actions and offer
rewards. Q asks when those tools make the agent report truthfully and act as intended. Its oversight
example uses a weaker monitor to choose rewards for a stronger actor, which knows the relevant state
exactly \cite{Q}.

P, \emph{When a Donor Wrist Cannot Locate a Rotating Held Object}, examines a handover in which the
held object is hidden from the receiver. It reads a proposal to use the donor arm's known wrist pose
in place of the hidden object pose alongside a demonstration that a hand can rotate an object within
its grasp. P observes that the object may then move relative to the wrist. It gives a conditional
limit on wrist based receiving, and proposes a measurement test. It does not report matched
histories of wrist, motor readings and actions with different object poses \cite{P}.

The peers asked whether a receiver could be guided by Q's reports or rewards when the observations
in P fail to identify which way the object faces. They separated two questions. Could the receiver
choose the right approach from its information? Could a reward based on measured outcomes induce a
desired physical action? The connection is conditional throughout: P neither reports a strategic
donor nor supplies a signal distribution for Q's model.

\section{What was done}

The peers first identified the possible ambiguity at handover. A fixed wrist pose need not fix
object orientation if the object rotates within the hand. They treated the wrist, motor encoder
readings and action history as the receiver's available record. They then asked what follows if two
orientations can have the same record and need incompatible receiving actions.

One peer turned this into a bound on success from the record alone. In parallel, the other took Q's
oversight menu and replaced its actor's exact view of the state with a coarse signal. The peers
checked both calculations against Q and P. They corrected an early overreach: one matched pair
supports a claim about that pair, but says nothing by itself about how often ambiguity occurs in
ordinary handovers.

Next, they considered incentives separately from inference. A finite model of measured outcomes
showed when an observation based reward can support a desired action, including a simple case in
which no such reward can work. Finally, they checked a structured signal channel for which adapting
Q's menu restores its intended action. These were mathematical checks of stipulated models, not
handover experiments.

\section{Results}

\subsection{A limit from the receiver's observations}

Let \(H\) be either of two equally likely object orientations, labelled \(0\) and \(1\). Let \(Y\)
be everything the receiver observes before acting, including the wrist, encoder and action history.
Write \(P_0\) and \(P_1\) for the distributions of \(Y\) under the two orientations. Suppose that no
receiving action succeeds for both. The peers found that a receiver using only \(Y\) has average
success probability at most
\[
\frac{1+\operatorname{TV}(P_0,P_1)}{2}.
\]
Here \(\operatorname{TV}(P_0,P_1)\), total variation distance, is the largest difference between the
probabilities that \(P_0\) and \(P_1\) assign to the same possible event. To succeed, the receiver
must in effect distinguish the two orientations. The displayed expression is the best possible
probability of doing so from \(Y\). A report or reward calculated only from \(Y\) cannot add
information. Ada gave the argument in \texttt{ada/observation\_quotient.tex}, and Emmy checked it in
\texttt{emmy/connection.md}.

If \(P_0=P_1\), the bound is \(1/2\). That applies to a two-trial experiment made from a matched
pair, or to a population whose two conditional observation laws really coincide. It does not follow
for routine handovers from one matched pair. A donor with extra information about object pose could
also send an informative report; that report would enlarge the receiver's information.

\subsection{A limit from what a reward can measure}

Ada also considered a finite set of feasible physical actions and a measurement that assigns each
action an observed outcome. A reward may depend on that outcome, but cannot distinguish actions with
the same outcome. For each private signal, Ada compared the desired action with possible deviations.
These comparisons form a graph of observed outcomes. An observation based reward supports the
desired actions exactly when every cycle in this graph has a nonpositive total gain from deviation.
Summing the reward inequalities around a cycle proves one direction; assigning rewards from maximum
path gains proves the other. The full statement and proof are in
\texttt{ada/observation\_quotient.tex}.

The immediate case is simple. If an actor prefers a different feasible physical action that produces
the same measured outcome, the reward is identical for both actions. It cannot deter that deviation.
This result assumes the actor can choose between the physical actions and that their measured
outcomes are fixed. P establishes neither a strategic donor nor such a hidden deviation on its
hardware.

\subsection{Q's oversight menu under a changed signal}

In Q's oversight example, the state \(\theta\) is uniform on \([0,1]\). A strong actor observes it
exactly and has a preference bias \(b\); a weak monitor with bias \(w\) chooses a reward from a
menu. A false report of \(b\) can drive the actor to an extreme action when the shifted state
crosses an endpoint. That threat makes truthful choice of the reward best for the monitor under Q's
assumptions \cite{Q}.

Emmy changed the actor's observation to a signal saying only whether \(\theta\) exceeds \(1/2\). The
actor's best estimate of \(\theta\) is then either \(1/4\) or \(3/4\). A small false report shifts
its action while keeping it inside \([0,1]\), so Q's extreme action threat does not arise. For an
actor bias of \(1/2\), a monitor bias of \(1/10\), and an allowed report of \(9/20\), the monitor's
expected squared loss falls from \(37/1200\) under truth to \(7/300\) under the false report. Emmy's
calculation is in \texttt{emmy/connection.md}; Ada checked the values. This shows a failure of Q's
\emph{unchanged} menu after changing its signal assumption. It does not contradict Q's result.

The peers also checked a limited repair. Suppose a signal reveals \(\theta\) with probability
\(p>0\), and otherwise gives fresh uniform noise. Its posterior mean is \((1-p)/2+ps\), where \(s\)
is the observed signal value. After rescaling actions, biases and rewards, Q's menu argument applies
to this signal and implements that posterior mean. The remaining expected squared error from
incomplete information is \((1-p^2)/12\). This channel was introduced for analysis; P gives no
evidence that handover observations follow it.

\section{Objections and remaining uncertainty}

The peers objected to treating these calculations as an established result about the paired systems.
P already points to pose aware handover work and proposes a sensor test. Q already recognises that
rewards depend on what can be measured. The finite graph condition is a general difference
constraint argument. The oversight calculation changes Q's exact observation premise, while its
repair relies on an invented signal channel. An adjacent haptic pose estimation study was recorded
during the prior work search, but was not used to prove a claim or establish novelty \cite{haptic}.

The cross checks settled two logical points. The original oversight menu can fail under a coarse
signal, yet a menu adapted to a particular signal can succeed. Likewise, identical histories in a
matched pair do not imply identical observation distributions in normal operation. What remains open
is physical: whether different object poses occur under matched histories, whether they demand
incompatible receiving actions, and whether an actor can choose hidden physical deviations. A common
safe grasp or negligible overlap would weaken the proposed connection.

\clearpage
\section{Why the thread stopped}

The verifier paused the thread because P does not show matched observation histories with different
object poses or show that those poses require incompatible receiving actions. The information bound
is therefore conditional. The oversight example changes Q's signal assumption and establishes no
handover result. The smallest step to restart the thread is to hold the donor wrist fixed,
externally measure a marked object's pose in occluded trials, and identify trials with matched
encoder and action histories. The receiving actions for any pose distinct matches must then be
tested for incompatibility. Only those measurements would show whether the mathematical limits bear
on this handover.

\bibliographystyle{plainurl}
\bibliography{references}

\end{document}
